\section{A pre-registered neural pilot that could not be resolved}\label{sec:pilot}
We pre-registered a neural study in the known-scale family (questions: population-like versus context-only behaviour in the
identifiable cells; library size at which memorisation gives way to generalisation), froze its code, configuration, decision
rules and seed table, and ran a pilot with dedicated seeds. The frozen rules chose the training budget as the smallest grid
point at which every control passed and the validation log loss at the shortest contexts stopped improving by more than
\ConvTol nats; they then capped the number of models per arm by a compute budget of \CapBudgetHours GPU-hours, with a minimum
of \MinModels models per arm, set before the pilot to give adequate power with a bootstrap-calibrated critical value.

The pilot chose \PilotB training steps: at the previous grid point the loss still improved by \PilotImprovementTwentyFive
nats, at this one by \PilotImprovementFifty. At that budget the cap allowed \PilotCap models per arm, below the minimum, so the
pre-registered outcome is \textsc{not resolvable (compute)}. This is the result of the pre-registered study. No effect-bearing
evaluation of the pilot models was run. A post-pilot budget increase was prepared, labelled as such, and not adopted; neural
effort was redirected to the unknown-noise family, whose map leaves more room (Section~\ref{sec:results}). We report this
outcome because a pilot that ends by its own rules is evidence about what the design can resolve, not a failure to hide.
